\subsection{Mean-aware representations}

The centring result suggests testing when the primary signs should
encode deviations from a mean (\cref{sec:8-centring}). Repeated cold
builds of centred and uncentred data should first separate this effect
from training variation. List-dependent means are another candidate,
but require consistent query scoring and residual reconstruction.
The test should compare shortlist coverage, the best attainable refined
recall, selected budgets and total throughput, so a cheaper shortlist
cannot conceal a lower quality ceiling.

\subsection{Calibration risk and reuse cost}

Can a statistically controlled budget choice save more query time than
it costs to obtain? Under the nesting conditions of
\cref{eq:6-rank-curve}, reference-neighbour primary ranks could replace
repeated trial searches. An earlier four-setting offline replay found
that sequential sampling added searches; a model fitted to measured
calibration times estimated higher cost despite fewer sampled queries.
It was not an integrated GPU controller measurement. A replacement
should test batched trials or rank reuse against a quality-matched fixed
budget using direct calibration-plus-reuse time. Its risk analysis must
cover the sampling and budget-selection procedure, and drift should
trigger renewal. Reference-relative control still needs a link to labels
to provide absolute recall guarantees.

\subsection{Residual storage across GPU and host memory}

Residuals account for 80.7--82.9\% of the total resident process memory
on BGE and Stella (\cref{sec:7-7}), but only shortlisted rows are read.
Keeping routing and primary codes on the GPU while fetching residual
blocks from host memory could increase capacity. The difficulty is
turning scattered shortlist positions into efficient block transfers.
Under fixed scoring and selection, vary the device budget and measure
extra bytes, cache hits, transfer overlap, throughput and tail latency.
This would establish the capacity--latency trade-off rather than assume
that offloading improves speed.

\subsection{Explaining the small-batch boundary}

At recall 0.97, grouping improves JHQ's small-batch execution yet does
not close the gap to CAGRA fp32 (\cref{sec:7-smallbatch}). Candidate
work, selection overhead and interface costs could each contribute.
A matched study should count candidates scored by both indexes and
separate kernel time from transfers and host-call overhead. Comparing
those costs across batches would test whether fewer probed candidates,
cheaper selection or a different index is the useful next step.
